\chapter{Metodologia}
\label{cap:metodologia}

Este capítulo descreve o caminho adotado para conduzir o trabalho e para avaliar se a solução proposta atende ao problema identificado. Seguindo uma organização \textit{top-down}, a metodologia apresenta primeiro a natureza da pesquisa, depois os procedimentos de desenvolvimento e, por fim, a estratégia de avaliação do protótipo.

\section{Caracterização da Pesquisa}
%DEGAS - NÃO SE TRATA DE UM TRABALHO DE PESQUISA. MAS DE DESENVOLVIMENTO.

%DEGAS "caracteriza-se como uma atividade de natureza tecnológica"

%DEGAS "O projeto possui abordagem qualitativa, 

%DEGAS - defina Top Down, e ponha no contexto no contexto. 

%DEGAS "metodologia apresenta primeiro a natureza da pesquisa necessária para compreensão do problema,"

O trabalho caracteriza-se como uma pesquisa aplicada de natureza tecnológica, pois tem como objetivo produzir uma solução computacional para apoiar um processo acadêmico real: a organização e a pré-validação do barema de atividades complementares do Colegiado de Ciência da Computação. A pesquisa possui abordagem qualitativa, uma vez que parte da análise do fluxo utilizado pelos discentes, das regras institucionais e das dificuldades observadas na preparação dos documentos.

O foco do estudo não é propor um novo método teórico de avaliação acadêmica, mas transformar regras já existentes em uma ferramenta funcional. Assim, o mérito do trabalho está na construção, organização e validação de um protótipo capaz de reduzir erros recorrentes no preenchimento, anexação e conferência inicial dos certificados.

\section{Procedimentos Metodológicos}

O desenvolvimento foi conduzido por prototipagem evolutiva, em que versões sucessivas do sistema foram refinadas a partir das necessidades identificadas durante o levantamento e a implementação \cite{soareseresende2007}. Essa estratégia foi adotada porque o problema envolve interação com usuário, regras de negócio específicas e tratamento de documentos em formato PDF, aspectos que exigem ajustes progressivos ao longo do desenvolvimento.

Inicialmente, foi realizado o levantamento das regras do barema e das informações necessárias para compor a solicitação do Em seguida, definiu-se a estrutura geral do sisdiscente. tema, incluindo a separação entre interface, rotas da aplicação, serviços de processamento, regras de validação e geração do documento final. Na etapa de implementação, foram desenvolvidos os mecanismos de preenchimento do formulário, anexação de certificados, extração de informações dos arquivos, pré-validação e montagem do PDF consolidado.

%DEGAS -  "Em seguida, definiu-se a estrutura geral do sisdiscente. tema, incluindo a separação entre interface, rotas da aplicação, serviços de processamento, regras de validação e geração do documento final" - Utilize um Itemize e defina rapidamente cada um. Na verdade, Este parágrafo acima traz muita informação junta. Ficou atropelado. Faz um explicando o levantamento das regras, outro falando da estrutura (com o itemize, talvez). Outro da implementação.

Durante o desenvolvimento, as funcionalidades foram testadas de forma incremental. A cada novo recurso, foram verificados os efeitos sobre o fluxo completo de uso, evitando que decisões de implementação ficassem desconectadas do objetivo central do sistema.

\section{Estratégia de Avaliação}

A avaliação do protótipo foi planejada com base em testes funcionais e em critérios de conformidade com as regras do barema. Foram considerados cenários como anexação de múltiplos certificados, identificação de carga horária, comparação entre horas solicitadas e horas comprovadas, detecção de certificados duplicados, geração do documento final e funcionamento da pré-validação

%DEGAS Mesma coisa. Explique como foi o planejamento dos testes e a avaliação de conformidade em um parágrafo. Mencione os cenários com detalhes em outr parágrafo, e use um Itemize pra separar cada um

\begin{itemize}
    \item Anexação de múltiplos certificados
    \item Identificação de carga horária
    \item Comparação de horas
    \item Detecção de certyificados duplicados
    \item Geração de dpcumento final
\end{itemize}

%DEGAS - "funcionamento da pré-validação" ficou complicado. O que é pré-validado? Talvez isso já esteja diluído na parte de cima do parágrafo (que vão virar múltiplos parágrafos). Mas veja aí. 

Também foi prevista a verificação do comportamento do sistema diante de limitações do processamento automático, como certificados sem texto extraível ou indisponibilidade de serviços externos de inteligência artificial. Nesses casos, a avaliação observa se o sistema mantém um comportamento seguro, recorrendo à validação básica quando necessário.

Dessa forma, a avaliação busca responder se o protótipo reduz erros antes da submissão oficial, mantém aderência às regras do COLCIC e produz um documento final organizado para análise posterior.
